% ============================================================================
\chapter{El cierre de cada mes}\label{cap:cierre}
% ============================================================================

El cierre se hace el \textbf{primer día hábil del mes siguiente} y lleva unos 30 minutos. Al terminarlo, los resultados del mes quedan definitivos.
El ejercicio~5 (sección~\ref{sec:ej5}) es un cierre completo de práctica.

\section{Lista de comprobación del cierre}

\begin{center}
\renewcommand{\arraystretch}{1.35}
\begin{tabularx}{\linewidth}{@{}c>{\raggedright\arraybackslash}X l@{}}
\toprule
& \textbf{Tarea} & \textbf{Hoja} \\ \midrule
$\square$ & 1. Todos los contenedores del mes están registrados y en \textbf{Cerrado}. & \hoja{ENTRADA} \\
$\square$ & 2. Las alertas del mes están resueltas o explicadas en Observaciones. & \hoja{ENTRADA} \\
$\square$ & 3. Los gastos fijos reales del mes que difieren del valor base están escritos en la columna del mes. & \hoja{COSTOS\_FIJOS} \\
$\square$ & 4. Si cambió algún precio durante el mes, está como una nueva fila con su fecha. & \hoja{PRECIOS} \\
$\square$ & 5. El mes a liquidar está elegido y la base variable de cada gestor, escrita. & \hoja{COMERCIAL} \\
$\square$ & 6. La conciliación dice \textbf{CUADRA}, o la diferencia está explicada (cobros pendientes). & \hoja{COMERCIAL} \\
$\square$ & 7. Los cuadres dicen \textbf{OK} y la columna «Control» de \hoja{RESUMEN\_MES} es 0. & \hoja{INICIO} \\
$\square$ & 8. Resultado del mes revisado: utilidad, costo por kg, punto de equilibrio, regiones. & \hoja{RESUMEN\_MES} \\
$\square$ & 9. Informe del mes exportado a PDF. & \hoja{DASHBOARD} \\
$\square$ & 10. Archivo guardado y copia de respaldo del mes hecha. & --- \\
\bottomrule
\end{tabularx}
\end{center}
El anexo~\ref{anx:checklist} tiene esta lista en una página aparte para imprimir.

\section{Detalle de cada tarea}

\subsection*{Gastos fijos reales (tarea 3)}
En \hoja{COSTOS\_FIJOS} cada mes copia el valor base. Compare con las facturas y nóminas del mes. Si un gasto fue distinto, escriba el importe real
en la columna de ese mes (los meses empiezan en la columna G, con el primer mes del período).

\begin{atencion}
\textbf{No cambie el «Valor base» (columna F) para corregir un solo mes}: cambiaría todos los meses, incluidos los ya cerrados.
Cambie el valor base solo si el gasto cambia \textbf{desde ahora y para siempre}, y antes siga el procedimiento del capítulo~\ref{cap:cambios}.
\end{atencion}

\subsection*{Liquidación de los gestores (tareas 5 y 6)}
\begin{pasos}
\paso En \hoja{COMERCIAL}, escriba en \celda{C4} el primer día del mes que cierra.
\paso En la columna \textbf{G} escriba lo que gestionó cada gestor en el mes: USD cobrados o kg, según la base.
\paso La columna \textbf{I} es lo que hay que pagar a cada uno y la fila 33, el total.
\paso Mire la conciliación (filas 36 a 39). Si dice REVISAR, compare gestor por gestor con los contenedores del mes.
      Las causas típicas son un cobro sin declarar, un contenedor sin registrar o un cobro pendiente.
\end{pasos}

\begin{nota}
Si la empresa paga la comisión sobre lo \textbf{efectivamente cobrado} (es lo recomendable), puede que al cierre falte cobrar algo.
En ese caso la diferencia de la conciliación es lo pendiente de cobro: anótelo y compruebe que se cobre el mes siguiente.
\end{nota}

\section{La copia de respaldo}\label{sec:respaldo}

Como usted es la única persona que lleva el libro, la copia de respaldo es su seguro: si el archivo se daña, se borra o se rompe la computadora,
el trabajo no se pierde.

\begin{pasos}
\paso Guarde el archivo maestro (\tecla{Ctrl}+\tecla{G}) y ciérrelo.
\paso Abra la carpeta del archivo en el Explorador de Windows. Haga clic derecho sobre el archivo $\rightarrow$ \menu{Copiar}, y luego clic derecho en
      un espacio vacío $\rightarrow$ \menu{Pegar}.
\paso Cambie el nombre de la copia (clic derecho $\rightarrow$ \menu{Cambiar nombre}) por \escribal{Respaldo\_\allowbreak 2026-10\_\allowbreak Rentabilidad.xlsx}
      (año y mes del cierre).
\paso Guarde también una copia \textbf{fuera de la computadora}: memoria USB, correo electrónico o la nube.
\end{pasos}

\begin{consejo}
Haga además una copia rápida cada viernes. Si un día comete un error que no puede deshacer, recupere la última copia y vuelva a escribir solo lo de
los últimos días. \textbf{Nunca trabaje sobre una copia de respaldo}: siempre sobre el maestro.
\end{consejo}
